\begin{abstract}
    % \lw{Paper title suggestions: Keypoint Trajectories for Functional Tool-Use Generalization} 
    % The purpose of this document is to provide both the basic paper template and submission guidelines. Abstracts should be a single paragraph, between 4--6 sentences long, ideally. Gross violations will trigger corrections at the camera-ready phase.
    % We define \textit{Functional Generalization} as the ability to accomplish the same function with diverse tools, such as using a hammer, a stone, or even a shoe to strike a nail. While humans naturally exhibit such flexibility in tool use, the capability of embodied agents to use tools for manipulation remains largely underexplored. 
    % Bridging this gap requires agents to capture the functional role of different tools beyond their visual appearance, which is challenging for two key reasons.
    % First, exposing agents to diverse tool-use behaviors typically requires large-scale robotic demonstrations, which are time-consuming and labor-intensive to collect. Second, directly learning sparse motor actions from dense visual observations often causes policies to overfit the discrete action patterns and overlook function-relevant information, limiting their ability to generalize to novel tools.  
    % In this paper, we propose the Generalizable Functional Plan and Execution (\algoname{}) policy, a System-2/System-1 framework, to enhance functional generalization in tool-use manipulation. Instead of directly mapping raw visual observations to motor commands, \algoname{} explores a functional intermediate representation that is visually predictable yet structured enough for action grounding. Specifically, the System-2 planner predicts the future states of the functional intermediate representation, while the System-1 execution policy translates the predicted states into executable actions. 
    % To evaluate functional generalization, we further introduce a hitting-function benchmark across seven tools. Extensive experiments show that our \algoname{} policy consistently improves generalization to novel tools compared to other state-of-the-art methods.    
    % Humans can wield a hammer, a stone, or even a shoe to drive a nail, because tools that share a function also share an \textit{affordance}: a common functional contact region and a common functional goal. 
    % Yet such affordance similarity, stable in visual and geometric space, does not transfer to raw action space, where each tool demands a distinct motor pattern, causing end-to-end policies to overfit to seen tools and fail on novel ones. 
    % We formalize this problem as \textit{Functional Generalization}: accomplishing the same function with diverse and unseen tools, and ask which prior best carries affordance from perception to action. 
    % Comparing affordance images, human video prompts, and 2D keypoint trajectories, we find that keypoint trajectories offer the best balance between affordance expressiveness and action groundability. 
    % Building on this finding, we propose \textbf{F}uncti\textbf{O}nal \textbf{R}easoning and \textbf{G}rounded \textbf{E}xecution (\algoname{}), a policy first predicts future keypoint trajectories from action-free observations and then grounds them into executable robot actions using a small set of action-labeled demonstrations. 
    % We further introduce a hitting-function benchmark across seven tools, where \algoname{} consistently outperforms state-of-the-art methods on unseen tools, achieving over $2\times$ improvement in average success rate.
While humans readily repurpose a book, a stone, or a shoe to drive a nail, robots trained on specific tools fail to transfer the same function to novel ones -- a gap we formalize as \textit{functional generalization}. Such tools share a common functional intent that is visually recognizable, yet this perceptual similarity does not carry over to action space, where each tool demands an entirely different motor pattern. To bridge this gap, we explore intermediate representations including affordance images, human video prompts, and 2D keypoint trajectories, finding that keypoint trajectories best balance functional expressiveness and action groundability. Building on this, we propose \textbf{F}uncti\textbf{O}nal \textbf{R}easoning and \textbf{G}rounded \textbf{E}xecution (\algoname{}), a two-stage policy that decouples functional reasoning from action execution: predicting generalizable keypoint trajectories from action-free data, then grounding them into robot actions with limited demonstrations. On a seven-tool hitting-function benchmark, \algoname{} consistently outperforms state-of-the-art methods on unseen tools in both simulation and the real world, achieving over $2\times$ improvement in average success rate. More experimental results are available at: \url{https://FORGE-policy.github.io}.
\end{abstract}

% Two or three meaningful keywords should be added here
% \keywords{Functional Generalization, \lw{Affordance, Robotic Tool Use, Keypoint Trajectories}}

\keywords{Affordance, Robotic Tool Use, Keypoint Trajectories}


% \lw{
% Humans can wield a hammer, a stone, or even a shoe to drive a nail, because tools that share a function also share an \textit{affordance}: a common functional contact region and a common functional goal. 
% Yet such affordance similarity, stable in visual and geometric space, does not transfer to raw action space, where each tool demands a distinct motor pattern, causing end-to-end policies to overfit to seen tools and fail on novel ones. 
% We formalize this problem as \textit{Functional Generalization}: accomplishing the same function with diverse and unseen tools, and ask which prior best carries affordance from perception to action. 
% Comparing affordance images, human video prompts, and 2D keypoint trajectories, we find that keypoint trajectories offer the best balance between affordance expressiveness and action groundability. 
% Building on this finding, we propose \algoname{} (Generalization via Functional Trajectories), which first predicts future keypoint trajectories from action-free observations and then grounds them into executable robot actions using a small set of action-labeled demonstrations. 
% We further introduce a hitting-function benchmark across seven tools, where \algoname{} consistently outperforms state-of-the-art baselines on unseen tools, achieving over $2\times$ improvement in average success rate.
% }
